%% icmi2026_preprocessing_temporal.tex
%% Target: 8-page ICMI/CSMI workshop paper on preprocessing and temporal modelling.
%%
%% Compile:
%%   pdflatex icmi2026_preprocessing_temporal
%%   bibtex   icmi2026_preprocessing_temporal
%%   pdflatex icmi2026_preprocessing_temporal
%%   pdflatex icmi2026_preprocessing_temporal

\documentclass[sigconf,review,anonymous]{acmart}

\usepackage{booktabs}
\usepackage{multirow}
\usepackage{amsmath}
\usepackage{enumitem}
\usepackage{xcolor}
\usepackage{microtype}

\acmConference[CSMI '26]{Workshop on Collective States in Multimodal Interaction (CSMI)}
  {October 2026}{Naples, Italy}
\acmYear{2026}
\copyrightyear{2026}
\settopmatter{printacmref=false}
\renewcommand\footnotetextcopyrightpermission[1]{}
\pagestyle{plain}

\definecolor{bestcol}{HTML}{1a6e1a}
\newcommand{\best}[1]{\textbf{\textcolor{bestcol}{#1}}}

\title{From Windows to Sequences: Preprocessing and Temporal Modeling Lessons
from GroupAffect-4}

\author{Anonymous Authors}
\affiliation{\institution{Anonymous Institution}}
\email{anon@example.com}

\begin{document}

\begin{abstract}
In small physiological affect datasets, preprocessing choices can dominate
model architecture.
This paper studies that effect in GroupAffect-4, a collaborative group dataset
with wearable physiology, eye tracking, and sparse Valence--Arousal--Dominance
self-reports.
Using a fixed train/validation/test split and training-only preprocessing
calibration, we compare 15s, 30s, 45s, and 60s temporal contexts; full-feature
and modality-selective smoothing; drop-one-modality ablations; enhanced
gaze/pupil feature extraction; optional speech feature inclusion; and
per-modality temporal encoders with SimCLR/MSM pre-training.
The strongest balanced preprocessing recipe smooths only EDA and PPG features
with centred 45s context while leaving gaze, pupil, and IMU unsmoothed,
improving all VAD dimensions over the unaugmented 15s baseline
($V=0.600$, $A=0.624$, $D=0.436$).
Reprocessing gaze and pupil features from raw sequences with improved NaN
handling changes the Dominance trade-off (up to $D=0.492$ in unsmoothed
all-modality v2), but under canonical centred-smoothing comparisons the
confirmed single-pipeline best remains $0.553$ mean macro-F1.
Adding speech to the v2 all-modality feature set does not improve this
benchmark and is therefore treated as a secondary analysis rather than a
default recipe.
A per-dimension canonical upper bound reaches $0.561$ and is reported as
exploratory.
Modality ablations show that IMU is the largest arousal contributor, followed
by EDA and PPG.
Temporal Conv1D and BiGRU models are trainable but do not beat the SVM
baseline as single shared configurations; however, SimCLR pre-training on
unlabelled sequences enables per-dimension configurations to exceed SVM+AP1 on
Arousal and Dominance.
The results motivate a practical recipe: use modality-aware preprocessing as
the default benchmark, and treat temporal self-supervision as a separate
sequence-modelling track requiring validation-only configuration selection.
\end{abstract}

\begin{CCSXML}
<ccs2012>
  <concept>
    <concept_id>10010147.10010257.10010293</concept_id>
    <concept_desc>Computing methodologies~Machine learning approaches</concept_desc>
    <concept_significance>500</concept_significance>
  </concept>
  <concept>
    <concept_id>10003120.10003121.10003129</concept_id>
    <concept_desc>Human-centered computing~Collaborative interaction</concept_desc>
    <concept_significance>300</concept_significance>
  </concept>
</ccs2012>
\end{CCSXML}

\ccsdesc[500]{Computing methodologies~Machine learning approaches}
\ccsdesc[300]{Human-centered computing~Collaborative interaction}

\keywords{physiological preprocessing, temporal modelling, wearable sensing,
SimCLR, EDA, PPG, IMU, affect recognition}

\maketitle

\section{Introduction}
\label{sec:intro}

Physiological affect recognition is often framed as a modelling problem:
given enough sensor streams, a sufficiently expressive temporal model should
learn affective state.
In sparse-label datasets, that framing is incomplete.
When only a few hundred labelled windows are available, the representation
presented to the model--window size, smoothing, feature aggregation, and split
isolation--can matter as much as the architecture itself.

This paper studies preprocessing and temporal modelling in the
GroupAffect-4 benchmark.
We ask two questions:
(1) what temporal context should physiological features use, and
(2) when do raw-sequence encoders add value beyond summary features?
We intentionally treat these as separate questions because they operate at
different levels of the pipeline.
Windowing and smoothing define what the supervised learner is allowed to see.
Temporal encoders define how that information is represented once the input
space has been fixed.
Conflating the two can lead to misleading conclusions: a neural sequence model
may appear weak simply because the 15s feature windows truncate slow EDA
responses, while an SVM may appear strong because handcrafted features already
encode the relevant temporal context.

The goal is therefore not to produce a single leaderboard winner.
Instead, we use GroupAffect-4 as a compact but realistic benchmark for small
group physiology: few labelled affect reports, many unlabelled sensor windows,
heterogeneous modalities, and strong temporal structure.
This setting is common in field and lab studies of collaboration, where
self-report must remain sparse to avoid interrupting interaction.

Our contributions are:
\begin{itemize}[leftmargin=1.2em,itemsep=1pt,topsep=2pt]
  \item We compare hard-pooled windows, full-feature smoothing, and
  modality-selective smoothing under a shared SVM benchmark.
  \item We show that smoothing slow EDA/PPG features while preserving fast
  gaze, pupil, and IMU features improves all VAD dimensions over the
  unaugmented 15s baseline while preserving sample count.
  \item We specify the AP1 pseudo-label augmentation protocol used in this
  paper so the recipe is self-contained and reproducible.
  \item We quantify modality contributions and show that IMU is the dominant
  arousal carrier in the current physiology-only feature set.
  \item We evaluate Conv1D and BiGRU temporal encoders, sequence augmentation,
  and SimCLR/MSM pre-training as an exploratory temporal modelling track.
\end{itemize}

\section{Related Work}
\label{sec:related}

\paragraph{Physiological affect recognition.}
Benchmark corpora including DEAP~\cite{koelstra2012deap},
WESAD~\cite{schmidt2018wesad}, DREAMER~\cite{katsigiannis2018dreamer}, and
MAHNOB-HCI~\cite{soleymani2012mahnob} established dimensional affect prediction
from EDA, cardiac, and wearable physiological signals.
Reviews confirm that EDA and cardiovascular responses are informative but
highly sensitive to preprocessing and splitting
choices~\cite{shu2018review,ahmad2022survey,kreibig2010autonomic}.
EDA phasic SCR responses unfold over 5--30 seconds, making window length a
first-order modelling choice rather than a tuning detail.
Pupil and gaze dynamics reflect cognitive load and arousal at faster
timescales~\cite{partala2003pupil,laeng2012pupillometry},
and body movement carries engagement and arousal signals in social
settings~\cite{karg2013body}.
These modality-specific timescale differences motivate this paper's central
hypothesis: a single window size and smoothing rule applied uniformly across
modalities is suboptimal.

\paragraph{Group and social affect datasets.}
K-EmoCon~\cite{park2020kemocon} and RECOLA~\cite{ringeval2013introducing}
contributed naturalistic dyadic interaction datasets with wearable sensors;
AMIGOS~\cite{mirandacorrea2021amigos} explicitly studied affect and personality
in group and individual conditions.
Meeting and social-signal corpora (AMI~\cite{mccowan2005ami},
SALSA~\cite{alameda2015salsa},
MatchNMingle~\cite{cabrera2018matchnmingle},
ELEA~\cite{sanchez2012emergent},
UDIVA~\cite{palmero2021context}) capture group behaviour, leadership, and
interaction dynamics.
GroupAffect-4~\cite{groupaffect4} adds per-person wearable physiology, Tobii
eye tracking, and post-task VAD self-reports to four-person structured
collaboration, occupying the intersection of physiological affect prediction
and group interaction research.

\paragraph{Self-supervised and semi-supervised learning for wearables.}
Time-series augmentation improves wearable classifiers at small label
counts~\cite{um2017data,iwana2021augmentation}.
Contrastive methods such as TS-TCC~\cite{eldele2021tscc} learn representations
from unlabelled physiological segments; SimCLR-style contrastive
pre-training~\cite{chen2020simclr} rewards invariance across augmented views of
the same window and has shown benefit in low-label regimes.
Semi-supervised approaches including Mean Teacher~\cite{tarvainen2017mean} and
FixMatch~\cite{sohn2020fixmatch} confirm that unlabelled data help only when
pseudo-label confidence is reliable~\cite{lee2013pseudo},
a finding directly relevant to our temporal pre-training track.
GroupAffect-4 provides a useful stress test: over 8\,000 unlabelled windows
but only 103 labelled training windows in the canonical split.

\paragraph{Personality and affect dynamics.}
The Big Five (BFI-44)~\cite{john1999bigfive} is the dominant trait
taxonomy for predicting affect.
ASCERTAIN~\cite{subramanian2018ascertain} and AMIGOS~\cite{mirandacorrea2021amigos}
show that personality shapes physiological affect responses.
\citet{meng2026moderating} demonstrated trait-specific moderation timescales
in VAD dynamics during real interaction, motivating per-dimension modality
conditioning.
Emotional contagion~\cite{hatfield1993emotional} provides the social
mechanism underlying cross-person label transfer in the AP1 protocol used
here.

\paragraph{Preprocessing as implicit modelling.}
In physiological affect recognition, preprocessing is often described as a
technical prelude to modelling.
For small datasets, it is more accurate to view preprocessing as part of the
model class.
A 15s mean EDA feature, a 45s centred EDA feature, and a raw 200-step EDA
sequence encode different assumptions about how long affective information
persists and how quickly labels should change.
Those assumptions can dominate the downstream learner when the number of
labelled windows is small.

\paragraph{Leakage and deployment time.}
Temporal smoothing also changes the evaluation question.
Centred smoothing uses both past and future neighbouring windows and is valid
for offline analysis, post-hoc annotation support, and dataset benchmarking.
It is not an online recipe for real-time intervention unless the future context
is replaced by a causal filter.
Throughout this paper, all scalers, thresholds, and class weights are fitted on
training windows only; when centred smoothing is used, it is described as an
offline preprocessing choice rather than an online claim.
This distinction is especially important for reproducible benchmarking because
the same score can imply very different deployment assumptions.

\section{Dataset and Features}
\label{sec:dataset}

GroupAffect-4~\cite{groupaffect4} contains 39 complete participants
(10 four-person groups) recorded across five sequential tasks:
\textbf{T0} (passive rest), \textbf{T1} (information pooling),
\textbf{T2} (negotiation), \textbf{T3} (ideation), and
\textbf{T4} (public-goods game).
After each task, participants rated affect using the 9-point
Self-Assessment Manikin (SAM)~\cite{bradley1994sam} on Valence (V),
Arousal (A), and Dominance (D).
Each participant also completed the BFI-44~\cite{john1999bigfive}
before the session.
T4 did not administer a Dominance probe; T4 Dominance is therefore
excluded from all metrics.

Physiological and oculomotor signals were recorded from two wearable platforms
(Table~\ref{tab:devices}).
EmotiBit~\cite{makowski2021neurokit2} captured EDA, PPG (Green channel, best
SNR), skin temperature, and 6-axis IMU at $\sim$25\,Hz.
EDA was decomposed via cvxEDA~\cite{greco2016cvxeda} into phasic (SCR) and tonic
(SCL) components; HR and HRV (RMSSD, SDNN) were re-derived from raw PPG
peaks~\cite{elgendi2013systolic} after correcting a firmware unit
error in the original dataset release
(HR was stored at $\times\!\tfrac{1}{200}$ scale).
Each participant wore a Tobii Pro Glasses 3 unit (40--60\,Hz) for gaze and
pupil signals; invalid samples were forward-filled and bilateral pupil diameter
averaged across valid samples.

\begin{table}[t]
\caption{Recording devices and feature counts. Summary dims = handcrafted
features used by SVM baselines; Seq.\ dims = raw channels for temporal encoders.}
\label{tab:devices}
\centering
\setlength{\tabcolsep}{3pt}
\small
\begin{tabular}{lp{1.5cm}ccc}
\toprule
\textbf{Modality} & \textbf{Device} & \textbf{Rate} &
  \textbf{Sum.} & \textbf{Seq.} \\
\midrule
Gaze/fixation & Tobii G3 & 40--60\,Hz & 6  & 9 \\
Pupil/blink   & Tobii G3 & 40--60\,Hz & 7  & 3 \\
EDA           & EmotiBit & $\sim$25\,Hz & 13 & 5 \\
PPG/HRV       & EmotiBit & $\sim$25\,Hz & 7  & 3 \\
IMU           & EmotiBit & $\sim$25\,Hz & 16 & 6 \\
\midrule
Total         & --       & --          & 49 & 26 \\
\bottomrule
\end{tabular}
\end{table}

We use training-fitted midpoint-tertile labels and per-dimension
inverse-frequency class weights (fitted on training data only).
The canonical split is T0+T1 train ($N{=}103$), T2 validation ($N{=}78$),
and T3 test ($N{=}65$).
All reported scores are macro-F1 over three affect classes per dimension.
We report per-dimension scores because each VAD dimension is expected to have
different physiological carriers and temporal scales.

\begin{table}[t]
\caption{Canonical split. Task order is temporal; test labels are always later
than training labels but belong to the same known teams.}
\label{tab:split}
\centering
\setlength{\tabcolsep}{4pt}
\small
\begin{tabular}{llcc}
\toprule
\textbf{Split} & \textbf{Tasks} & \textbf{$N$} & \textbf{Role} \\
\midrule
Train      & T0 + T1 & 103 & Fit models, thresholds, scalers \\
Validation & T2      &  78 & Select configurations \\
Test       & T3      &  65 & Final held-out evaluation \\
\bottomrule
\end{tabular}
\end{table}

\paragraph{AP1 augmentation protocol used in this paper.}
AP1 denotes a fixed pseudo-labelled pool generated from training-task data
using Gaussian-process (GP) posteriors and personality-based trust weights
\cite{rasmussen2006gaussian,lee2013pseudo,john1999bigfive}.
For each VAD dimension, GP posteriors are mapped to Low/Mid/High classes using
training-fitted thresholds only.
Accepted unlabelled windows are merged with labelled training windows as
\begin{equation}
  \mathcal{D}_{aug}=\mathcal{D}_{L}\cup\mathcal{D}_{U},
\end{equation}
where each pseudo-labelled sample has weight
$\alpha_u=\mathrm{BFI\_sim}_u\in[0,1]$.
For participant $i$ in group $G$, the trust weight is
\begin{equation}
  \mathrm{BFI\_sim}_i=
  \frac{1}{|G|-1}\sum_{j\neq i}
  \frac{\mathbf{b}_i\cdot\mathbf{b}_j}{\|\mathbf{b}_i\|\|\mathbf{b}_j\|},
\end{equation}
where $\mathbf{b}_i$ is the BFI trait vector.
No T3 labels are used for pseudo-label generation, threshold fitting, scaler
fitting, or model selection.
In this paper AP1 is used as a fixed augmentation input, not tuned on the T3
test task.

\paragraph{Sequence representation.}
For temporal models, each modality is resampled independently to
$T=200$ timepoints per 15s window by piecewise linear interpolation.
This gives an effective temporal grid of approximately 13.3 Hz, sufficient for
the low-frequency physiological bands considered here.
Windows with too few valid raw samples in a modality are zero-filled for that
modality rather than interpolated from unreliable data.
The resulting tensor is represented as five modality-specific arrays:
\begin{equation}
  \mathbf{X} =
  \{\mathbf{X}^{gaze}_{T\times 9},
    \mathbf{X}^{pupil}_{T\times 3},
    \mathbf{X}^{EDA}_{T\times 5},
    \mathbf{X}^{PPG}_{T\times 3},
    \mathbf{X}^{IMU}_{T\times 6}\}.
\end{equation}
Sequence scalers fit a global mean and standard deviation on training windows
only, concatenating over time within each modality.
Validation, test, and unlabelled pool windows reuse those training parameters.

\section{Windowing and Smoothing}
\label{sec:windowing}

We evaluate three preprocessing families.
\textbf{Hard pooling} merges neighbouring 15s windows into longer contexts,
reducing the number of training samples.
\textbf{Full-feature smoothing} averages all feature groups with temporal
neighbours while preserving the original sample count.
\textbf{Physio-selective smoothing} averages only EDA and PPG features, leaving
gaze, pupil, and IMU untouched.
The distinction between pooling and smoothing matters: pooling changes both the
effective context and the number of supervised examples, while smoothing changes
the feature estimate but keeps the same label count.
At this scale, losing labelled examples can be as damaging as improving signal
quality.

\paragraph{Hard pooling.}
The 30s and 45s pooled settings merge consecutive 15s windows before feature
aggregation.
This is useful for slow autonomic responses but reduces the number of labelled
training examples from 103 to 74 and 57, respectively.
The 60s exact setting uses longer extraction context without the same
neighbour-merge rule and is included as a dilution check.
Formally, for adjacent 15s feature vectors $\mathbf{x}_{i}$ and
$\mathbf{x}_{i+1}$ within the same participant and task, the pooled 30s feature
is computed from the raw signal over the union interval rather than by simply
averaging the two feature vectors when a raw-signal extractor is available.
This matters for nonlinear features such as SCR peak count and HRV dispersion.
The price is that the number of supervised windows decreases.

\paragraph{Full-feature smoothing.}
Full smoothing averages each window's feature vector with nearby windows.
Forward smoothing is compatible with causal deployment, while centred smoothing
uses both previous and next windows and is therefore an offline recipe.
The hypothesis is that more stable estimates of noisy physiological features
should improve Arousal and Dominance, but this may blur fast oculomotor and
movement signals.
For a feature $f_{i}$ at window $i$, centred 45s smoothing uses
\begin{equation}
  \tilde{f}_{i} = \frac{1}{3}\left(f_{i-1}+f_i+f_{i+1}\right),
\end{equation}
within the same participant and task, with boundary windows using the available
neighbours only.
Forward smoothing replaces $f_{i+1}$ with the next available past window and is
therefore causal.

\paragraph{Physio-selective smoothing.}
The modality-selective variant applies the same temporal averaging only to EDA
and PPG/HRV features.
Gaze, pupil, and IMU features remain the original 15s estimates.
This implements a simple prior: slow autonomic signals benefit from a longer
effective context, while fast behavioural and oculomotor signals should retain
their local variation.

Table~\ref{tab:window} shows that there is no universal best window.
Hard-pooled 30s windows maximise Arousal but reduce Valence and training $N$.
Full-feature smoothing improves Arousal but damages Valence.
Physio-selective smoothing resolves the conflict: centred 45s EDA+PPG
smoothing with AP1 augmentation improves all three VAD dimensions while
retaining $N=103$.

\begin{table}[t]
\caption{Window size and smoothing ablation with RBF-SVM.}
\label{tab:window}
\centering
\setlength{\tabcolsep}{3.5pt}
\small
\begin{tabular}{llcccc}
\toprule
\textbf{Window} & \textbf{Aug} & \textbf{$N$} & \textbf{V} & \textbf{A} & \textbf{D} \\
\midrule
15s default       & none & 103 & 0.570 & 0.530 & 0.385 \\
15s               & AP1  & 103 & \best{0.608} & 0.535 & 0.435 \\
30s pooled        & none & 74  & 0.477 & \best{0.662} & \best{0.452} \\
45s pooled        & none & 57  & 0.468 & 0.404 & 0.413 \\
60s exact         & none & 103 & 0.458 & 0.519 & 0.364 \\
\midrule
Smooth-45s all    & AP1  & 103 & 0.517 & 0.640 & 0.434 \\
Physio-30s        & AP1  & 103 & 0.586 & 0.599 & 0.435 \\
\best{Physio-45s} & AP1  & 103 & 0.600 & 0.624 & 0.436 \\
\bottomrule
\end{tabular}
\end{table}

\paragraph{Interpretation.}
The best result depends on the target dimension.
If the goal is maximum Arousal alone, 30s hard pooling is strongest
($A=0.662$), consistent with SCR and pulse features needing more than one 15s
segment.
If the goal is maximum Valence, the original 15s AP1 recipe remains strongest
($V=0.608$), suggesting that local gaze, pupil, and movement statistics already
capture much of the relevant valence variation.
For a balanced devkit benchmark, however, physio-selective 45s smoothing is the
most defensible default because it improves all dimensions over the unaugmented
15s baseline while preserving the full labelled sample count.

\section{Modality Contributions}
\label{sec:modality}

Drop-one-modality ablations (Table~\ref{tab:modality}) show that IMU is the
largest contributor in the physiology-only feature set, driven by Arousal.
EDA and PPG contribute comparably; pupil and gaze are less important after
other autonomic and motion features are present.

\begin{table}[t]
\caption{Drop-one modality ablation, no augmentation. $\Delta$ is mean-F1 drop
relative to the full model.}
\label{tab:modality}
\centering
\setlength{\tabcolsep}{4pt}
\small
\begin{tabular}{lcccc}
\toprule
\textbf{Drop} & \textbf{V} & \textbf{A} & \textbf{D} & $\Delta$ \\
\midrule
None full  & 0.414 & 0.481 & 0.516 & -- \\
Gaze       & 0.489 & 0.477 & 0.377 & $-0.023$ \\
Pupil      & 0.464 & 0.479 & 0.360 & $-0.036$ \\
EDA        & 0.402 & 0.435 & 0.296 & $-0.093$ \\
PPG        & 0.396 & 0.411 & 0.323 & $-0.094$ \\
\best{IMU} & 0.431 & 0.269 & 0.361 & $\mathbf{-0.117}$ \\
\bottomrule
\end{tabular}
\end{table}

This pattern is consistent with collaborative tasks: wrist acceleration and
gyroscope features capture movement bursts, gestures, and engagement-related
motor activity.
The result does not imply that pupil or gaze are uninformative in general; it
shows that their marginal contribution is smaller once EDA, PPG, and IMU are
available in this dataset.
It also helps explain the smoothing results.
The largest Arousal gains come from stabilising slow autonomic features, but
the largest modality drop is IMU because movement is a strong marker of group
engagement in T3.
Smoothing IMU would therefore remove precisely the short bursts that make it
useful.
The preprocessing recommendation is consequently not ``smooth physiology''
in a broad sense; it is specifically to smooth EDA and PPG while preserving
movement features.


\section{Enhanced Gaze and Pupil Features}
\label{sec:v2_features}

The original 49-dimensional feature set uses basic summary statistics for gaze
(6 features: mean/std of $x$, $y$, validity rate, saccade proxy) and pupil
(7 features: mean/std of left, right, average diameter, blink rate proxy).
These are adequate for a baseline, but they discard the rich temporal dynamics
available in the raw Tobii G3 sequences.

We recompute gaze and pupil features from the stored 200-point per-window
sequences using signal-processing methods appropriate for the effective 13.3~Hz
sampling rate.
The enhanced gaze features (22 total) include fixation count, duration
statistics, dispersion, saccade count and amplitude, scanpath length and convex
hull area, spatial entropy, transition rate, and quadrant dwell proportions.
The enhanced pupil features (18 total) include blink rate and duration,
inter-blink intervals, LHIPA (Low/High-frequency Index of Pupillary Activity),
pupil derivatives, RMSSD, slope, range, skewness, kurtosis, and median.

\paragraph{Feature selection under the curse of dimensionality.}
We evaluate both full all-modality and reduced v2 feature sets.
The full all-modality representation has 82 features
(22 gaze $+$ 18 pupil $+$ 36 EDA/PPG/IMU $+$ 6 speech), while the
physio-only variant has 76 (same set without speech).
For augmentation-compatible comparisons, we additionally use the 47-feature
intersection shared between the v2 dataset and the AP1 pool.
A data-lineage audit explains the previously inconsistent train count: the v2
file was generated from the speech-enriched dataset file
(296 windows, T0+T1 train $=105$), whereas the canonical benchmark uses
\texttt{dataset\_15s.pkl} (292 windows, T0+T1 train $=103$).
The extra four windows (two in train) were inherited from the speech-enriched
pickle lineage and were not an intended protocol change.

\begin{table}[t]
\caption{Effect of enhanced gaze/pupil reprocessing on Dominance.
All rows use centred EDA/PPG smoothing on the T3 test split;
the all-modality row includes speech features (82 total).}
\label{tab:v2_dominance}
\centering
\setlength{\tabcolsep}{4pt}
\small
\begin{tabular}{p{0.42\columnwidth}cccc}
\toprule
\textbf{Feature set} & \textbf{V} & \textbf{A} & \textbf{D} & \textbf{Mean} \\
\midrule
Orig 49 + AP1 (smooth)           & 0.600 & 0.624 & 0.436 & 0.553 \\
V2 shared-47 (smooth, no AP1)    & 0.552 & 0.568 & 0.435 & 0.518 \\
V2 shared-47 + AP1 (smooth)      & 0.526 & \best{0.624} & \best{0.460} & 0.537 \\
V2 all-modal 82 (smooth, no AP1) & 0.460 & 0.514 & 0.416 & 0.463 \\
\bottomrule
\end{tabular}
\end{table}

Table~\ref{tab:v2_dominance} reveals that the v2 reprocessing dramatically
changes the VAD trade-off under the canonical split.
The shared-47 smooth setting does not improve Dominance over the original
smooth+AP1 baseline (0.435 vs.
0.436), while shared-47+AP1 recovers Dominance to 0.460 at the expense of
Valence.
The improvement comes not from new feature \emph{types} but from cleaner
feature \emph{values}: recomputing basic statistics (gaze\_x\_mean,
pupil\_avg\_std, etc.) directly from the stored sequences with proper NaN
interpolation produces more reliable estimates.
The all-modality 82-feature row confirms modality coverage but does not exceed
the shared-47 Dominance result, consistent with small-sample variance and
cross-modality dilution at this dataset size.

\paragraph{Speech inclusion decision.}
Speech is retained as an optional modality analysis, but it is not promoted to
the default benchmark recipe.
Under matched centred-smoothing comparisons in
Table~\ref{tab:v2_dominance}, the all-modality 82-feature setting (with speech)
does not outperform the shared-47 or original-49 recipes in mean macro-F1.
For this reason, all headline recipes in this paper remain physiology-first,
with speech reported as secondary evidence.

For benchmark comparability, all headline claims in this paper are anchored
to the canonical T0+T1/T2/T3 protocol (train $=103$).

AP1 effects on v2 rows are sensitive to preprocessing compatibility because
the augmentation pool was generated with the original pipeline.
This distributional mismatch can dilute gains and changes sign across settings,
especially for Dominance.
This motivates the per-dimension recipe below.

\subsection{Per-Dimension Optimized Recipe}
\label{sec:perdim_recipe}

The preprocessing and augmentation findings have complementary strengths
across VAD dimensions.
Table~\ref{tab:perdim_recipe} reports both a confirmed single-pipeline result
and a per-dimension upper bound.

\begin{table}[t]
\caption{Per-dimension optimized preprocessing recipe. Each dimension uses its
best-matched combination of feature processing and augmentation.}
\label{tab:perdim_recipe}
\centering
\setlength{\tabcolsep}{4pt}
\small
\begin{tabular}{llcccc}
\toprule
\textbf{Dimension} & \textbf{Recipe} & \textbf{V} & \textbf{A} & \textbf{D} & \textbf{Mean} \\
\midrule
Valence & Orig + smooth + AP1 & \best{0.600} & -- & -- & -- \\
Arousal & Orig + smooth + AP1 & -- & \best{0.624} & -- & -- \\
Dominance & V2 + smooth + AP1 & -- & -- & \best{0.460} & -- \\
\midrule
\multicolumn{2}{l}{Per-dim upper bound (exploratory)} & 0.600 & 0.624 & 0.460 & 0.561 \\
\multicolumn{2}{l}{\best{Confirmed single-pipeline (R2a)}} & \best{0.600} & \best{0.624} & 0.436 & \best{0.553} \\
\bottomrule
\end{tabular}
\end{table}

The per-dimension row should be interpreted as an exploratory upper bound.
Each component still follows a mechanistic hypothesis:
\begin{itemize}[leftmargin=1.2em,itemsep=1pt,topsep=2pt]
  \item \textbf{Valence/Arousal}: These dimensions benefit from AP1 because
  personality similarity carries information about affective appraisal
  while preserving training-task split isolation. The pool features match the
  original preprocessing.
  \item \textbf{Dominance}: This dimension benefits from cleaner gaze/pupil
  feature extraction. Under centred smoothing, v2+AP1 recovers Dominance to
  0.460, but this still underperforms the strongest unsmoothed v2-D setting;
  this sensitivity indicates a distribution-matching issue between pool
  generation and reprocessed features.
\end{itemize}

It demonstrates that preprocessing and augmentation interact differently
per dimension, but confirmatory reporting should use the single-pipeline
R2a value unless per-dimension selection is pre-registered on validation.

\section{Temporal Encoders}
\label{sec:temporal}

We evaluate per-modality Conv1D and BiGRU encoders over resampled raw
sequences.
The original two-track implementation zero-padded temporal inputs for
unlabelled pool windows; later inspection showed that pool sequences are
available, enabling corrected v2 training and self-supervised pre-training.
We keep both generations as useful baselines.

\paragraph{Architectures.}
The Conv1D encoder uses two strided temporal convolution layers per modality,
followed by global average pooling to produce a 32-dimensional embedding.
The BiGRU encoder uses a bidirectional recurrent layer per modality and mean
pools hidden states to produce a 64-dimensional embedding.
In both cases, modality embeddings are concatenated with the 49-dimensional
summary vector before the VAD prediction heads.
This hybrid design lets the model use raw temporal traces without discarding
the robust handcrafted statistics used by the SVM baseline.

\paragraph{Sequence augmentation.}
For labelled windows, v2 training applies lightweight time-domain
augmentation: Gaussian noise, amplitude jitter, and circular time shifts.
These perturbations are intentionally conservative because physiological
signals are not image-like; large time warps or amplitude changes can destroy
the relation between window and self-report.
The aim is regularisation, not synthetic label expansion.

\paragraph{Two uses of the unlabelled pool.}
The unlabelled pool can enter the pipeline in two different ways.
Label-space augmentation assigns GP-derived soft labels to pool windows and
mixes them into supervised training.
Representation pre-training uses pool sequences without labels, then fine-tunes
on the 103 labelled training windows.
The temporal experiments show that these roles should not be treated as
interchangeable: the pool is often more valuable for learning sequence
invariances than for adding many noisy supervised examples.

\paragraph{Pre-training objectives.}
The SimCLR condition samples two augmented views of the same multimodal window
and trains the encoder with an NT-Xent contrastive loss.
The positive pair is the same window under different perturbations; negatives
are other windows in the batch.
The MSM condition masks contiguous temporal chunks and reconstructs the held-out
values from the remaining context.
SimCLR therefore rewards invariance to mild temporal distortions, whereas MSM
rewards local predictive reconstruction.
The difference matters for affect: slow dominance-like state information may be
better preserved by contrastive identity, while short phasic dynamics may be
well served by reconstruction.

No single temporal configuration beats SVM+AP1 in mean F1
(Table~\ref{tab:temporal}), but corrected sequence use and augmentation close
much of the gap.
GRU benefits strongly from SimCLR pre-training on the full unlabelled pool in
the exploratory v3 run, especially for Dominance.

\begin{table}[t]
\caption{Temporal encoder comparison on the T3 holdout.}
\label{tab:temporal}
\centering
\setlength{\tabcolsep}{4pt}
\small
\begin{tabular}{lccc}
\toprule
\textbf{Model} & \textbf{V} & \textbf{A} & \textbf{D} \\
\midrule
MLP summary, no aug       & 0.568 & 0.505 & 0.375 \\
Conv1D v1, no aug         & 0.505 & 0.477 & 0.356 \\
GRU v1, no aug            & 0.449 & \best{0.592} & 0.385 \\
\midrule
Conv1D v2 + seqaug        & 0.511 & 0.513 & 0.429 \\
Conv1D v2 + seqaug + AP1  & 0.445 & 0.526 & 0.395 \\
GRU v2 + seqaug + AP1     & 0.459 & 0.500 & 0.454 \\
\midrule
Conv1D + SimCLR           & 0.601 & 0.506 & 0.351 \\
GRU + SimCLR              & 0.460 & 0.536 & \best{0.520} \\
\midrule
SVM + AP1 reference       & \best{0.608} & 0.535 & 0.435 \\
\bottomrule
\end{tabular}
\end{table}

\section{Pre-training Factorial}
\label{sec:factorial}

To separate architecture, pre-training, and fine-tuning augmentation effects,
we ran a $2\times3\times3$ factorial over encoder (Conv1D/GRU), pre-training
(none, SimCLR, MSM), and fine-tuning augmentation (none, default, time-warp).
The stricter factorial uses only the 3,127 windows from training tasks for
pre-training.

Table~\ref{tab:factorial} reports the best per-dimension configurations.
SimCLR is the strongest pre-training method on average, with the gain
concentrated in Dominance.
The important practical result is not that one temporal model wins globally;
it is that each VAD dimension prefers a different temporal recipe.

\begin{table}[t]
\caption{Factor-level marginal means from the temporal factorial. Each entry
averages over the other factors and three random seeds.}
\label{tab:factor_marginals}
\centering
\setlength{\tabcolsep}{4pt}
\small
\begin{tabular}{llcccc}
\toprule
\textbf{Factor} & \textbf{Level} & \textbf{V} & \textbf{A} & \textbf{D} & \textbf{Mean} \\
\midrule
\multirow{2}{*}{Encoder}
  & Conv1D & \best{0.510} & \best{0.521} & 0.383 & 0.471 \\
  & GRU    & 0.499 & 0.509 & \best{0.410} & \best{0.473} \\
\midrule
\multirow{3}{*}{Pre-train}
  & None   & 0.499 & 0.513 & 0.391 & 0.468 \\
  & SimCLR & 0.503 & 0.511 & \best{0.422} & \best{0.478} \\
  & MSM    & \best{0.512} & \best{0.521} & 0.376 & 0.470 \\
\midrule
\multirow{3}{*}{Aug.}
  & None    & 0.492 & \best{0.519} & \best{0.404} & 0.472 \\
  & Default & 0.500 & 0.517 & 0.394 & 0.470 \\
  & Time-warp & \best{0.521} & 0.509 & 0.392 & \best{0.474} \\
\bottomrule
\end{tabular}
\end{table}

\begin{table}[t]
\caption{Best per-dimension configurations from the temporal factorial.}
\label{tab:factorial}
\centering
\setlength{\tabcolsep}{3.5pt}
\small
\begin{tabular}{llccc}
\toprule
\textbf{Dim} & \textbf{Config} & \textbf{V} & \textbf{A} & \textbf{D} \\
\midrule
Valence & Conv1D + SimCLR + TW      & 0.553 & 0.544 & 0.378 \\
Arousal & Conv1D + SimCLR + default & 0.524 & \best{0.545} & 0.394 \\
Dominance & GRU + SimCLR + none     & 0.495 & 0.489 & \best{0.455} \\
\midrule
SVM + AP1 & reference               & \best{0.608} & 0.535 & 0.435 \\
\bottomrule
\end{tabular}
\end{table}

\paragraph{Main effects.}
Table~\ref{tab:factor_marginals} shows that encoder choice has little effect
on mean F1, but the dimensional pattern is meaningful.
Conv1D is stronger for Valence and Arousal, where short local transients and
phasic responses matter.
GRU is stronger for Dominance, where sustained interpersonal control and
engagement may unfold over longer intervals.
SimCLR is the best pre-training method overall, again driven by Dominance.
Masked sequence modelling (MSM) is competitive for Valence and Arousal but
lags on Dominance, suggesting that reconstructing local chunks is less useful
for slow state-like signals than learning invariance across augmented views.

\paragraph{Configuration selection.}
The per-dimension winners in Table~\ref{tab:factorial} exceed the SVM+AP1
reference on Arousal and Dominance, but they are selected from a factorial
grid.
For a final temporal-modelling submission, the selection rule should be locked
using validation folds before reporting T3 test results.
The present paper treats these results as evidence that self-supervised
sequence modelling is promising, not as a completed claim that temporal neural
models generally outperform the SVM baseline.

\section{Confirmed Benchmarks and Exploratory Tracks}
\label{sec:recipes}

The benchmark use case requires recipes that other researchers can reproduce.
Table~\ref{tab:recipes} summarises the benchmark settings we recommend for
future GroupAffect-4 experiments.
The first recipe is the conservative baseline: 15s summary features with no
augmentation.
The second is the fixed known-team AP1 augmentation baseline.
The third is the main single-pipeline preprocessing recipe.
The final two rows are explicitly exploratory and should not be reported as
drop-in replacements for confirmed baselines.

\begin{table}[t]
\caption{Benchmark recipe status map. Confirmed rows are suitable for direct
comparison in current workshop submissions; exploratory rows require
validation-locked confirmation.}
\label{tab:recipes}
\centering
\setlength{\tabcolsep}{2pt}
\footnotesize
\begin{tabular}{p{0.13\columnwidth}p{0.08\columnwidth}p{0.22\columnwidth}p{0.50\columnwidth}}
\toprule
\textbf{Status} & \textbf{Recipe} & \textbf{Use when} & \textbf{Core setting} \\
\midrule
Confirmed & R0 & Conservative baseline & 15s SVM, no augmentation \\
Confirmed & R1 & Known-team augmentation baseline & 15s SVM + AP1 \\
Confirmed & R2a & Balanced single-pipeline & EDA/PPG 45s smooth + AP1 ($V{=}0.600$, $A{=}0.624$, $D{=}0.436$) \\
Exploratory & E1 & Per-dim upper bound & R2a for V/A + v2+smooth+AP1 for D (mean $=0.561$) \\
Exploratory & E2 & Temporal representation & SimCLR pre-train + val-selected encoder \\
\bottomrule
\end{tabular}
\end{table}

R0, R1, and R2a are stable baselines for the current benchmark.
R2a is the main single-pipeline recommendation of this paper: it has a clear
mechanistic rationale and improves all three VAD dimensions relative to the
unaugmented 15s baseline.
E1 is an exploratory per-dimension upper bound and should be reported only
with explicit validation-locking caveats.
E2 is a research track: the best temporal configurations are encouraging, but
the final publication version should rerun configuration selection under a
strict validation-only protocol.

\section{Dimension-Level Mechanisms}
\label{sec:dimension_mechanisms}

The preprocessing and temporal results are not uniform across VAD dimensions.
This is expected: Valence, Arousal, and Dominance are not three interchangeable
classification heads.
They have different psychological meanings, different physiological carriers,
and different temporal scales.

\paragraph{Valence.}
Valence is the dimension least helped by long-window pooling.
The best Valence score remains the 15s AP1 SVM baseline, and full-feature
smoothing damages Valence sharply.
This suggests that Valence-relevant information in the current feature set is
local: gaze stability, pupil dynamics, facial-orienting proxies, and short
movement changes may matter more than slow autonomic averages.
Temporal encoders also struggle to beat the summary baseline for Valence.
One interpretation is that the handcrafted summary vector already captures the
useful local statistics at the current label density.

\paragraph{Arousal.}
Arousal is the clearest beneficiary of longer physiological context.
Hard-pooled 30s windows produce the highest Arousal score, and both full
smoothing and physio-selective smoothing improve Arousal over the default 15s
setup.
This is consistent with EDA phasic responses and PPG-derived HRV estimates,
which are unstable in very short windows.
The important caveat is that simply increasing context can harm other
dimensions.
The recommended recipe therefore smooths EDA and PPG but preserves faster
gaze, pupil, and IMU statistics.

\paragraph{Dominance.}
Dominance benefits from physio-selective smoothing and GRU+SimCLR pre-training.
This pattern suggests a slowly-varying state tied to sustained interpersonal
control and confidence.
However, Dominance is also the dimension most likely to need relational
features absent from the current physiology-only set.
Turn-taking, speaking time, gaze-to-person, posture, and role-specific task
actions may be necessary to move beyond the current ceiling.

\paragraph{Why one shared temporal model is insufficient.}
The factorial results show that each dimension prefers a different temporal
recipe.
This should not be treated as inconvenient noise; it is a substantive result.
If Valence is driven by local oculomotor and movement summaries, Arousal by
phasic autonomic context, and Dominance by sustained interpersonal state, then
forcing a single encoder/pre-training/augmentation recipe across all three is
biologically and psychologically implausible.
Future temporal models should either select configurations per dimension on
validation data or use architectures that explicitly represent dimension-level
timescales.

\section{Reproducibility Checks}
\label{sec:checks}

Small physiological datasets are vulnerable to subtle protocol drift.
Table~\ref{tab:checks} lists the checks we recommend when reusing these
benchmarks.
They are included because several apparent improvements in early devkit
iterations were caused by mismatched thresholds, pool handling, or split
semantics rather than genuine model improvements.

\begin{table}[t]
\caption{Reproducibility checks for preprocessing and temporal experiments.}
\label{tab:checks}
\centering
\setlength{\tabcolsep}{2pt}
\small
\begin{tabular}{p{0.33\columnwidth}p{0.59\columnwidth}}
\toprule
\textbf{Check} & \textbf{Reason} \\
\midrule
Fit scalers on train only & Prevents validation/test distribution leakage. \\
Fit label thresholds on train only & Keeps macro-F1 comparable across models. \\
Mark centred smoothing offline & Avoids claiming online performance from future context. \\
Keep T4 Dominance masked & T4 Dominance is not an administered label. \\
Separate pool pre-training from pseudo-labelling & Prevents noisy labels from overwhelming $N=103$. \\
Select temporal configs on validation & Avoids grid-searching directly on T3. \\
\bottomrule
\end{tabular}
\end{table}

\paragraph{Recommended reporting.}
Future papers should report the exact window recipe, whether smoothing is
causal or centred, whether AP1/AP2 augmentation is used, and whether unlabelled
windows are used for supervised pseudo-label training or self-supervised
pre-training.
They should also report per-dimension VAD scores, because a single mean can
hide the main scientific trade-off: Arousal often benefits from longer
physiological context, while Valence is more sensitive to preserving fast
signals.

\begin{table}[t]
\caption{Claim-status map for preprocessing and temporal modelling.}
\label{tab:claim_status}
\centering
\setlength{\tabcolsep}{3pt}
\small
\begin{tabular}{p{0.31\columnwidth}p{0.57\columnwidth}}
\toprule
\textbf{Claim} & \textbf{Current status} \\
\midrule
Physio-selective smoothing helps & Mature; improves all VAD dimensions over 15s no-aug. \\
30s pooling maximises Arousal & Supported, but trades off Valence and sample count. \\
IMU is largest arousal contributor & Supported by drop-one ablation, not a standalone sensor ranking. \\
Temporal pre-training can beat SVM per dimension & Promising; requires validation-locked confirmation. \\
One temporal config solves all VAD & Not supported; dimensions prefer different recipes. \\
\bottomrule
\end{tabular}
\end{table}

\section{Uncertainty and Evaluation Scope}
\label{sec:uncertainty_scope}

The claims are intentionally split by evidential strength.
R0, R1, and R2a are confirmatory under the fixed known-team split;
E1 and E2 remain exploratory and require validation-locked confirmation.
Centred smoothing is an offline recipe and should not be interpreted as a
real-time method; the same hypothesis should be retested with causal smoothing
for online deployment.

\section{Discussion}
\label{sec:discussion}

\paragraph{Preprocessing is a model choice.}
At $N=103$, feature smoothing and window definition change the benchmark as
much as model family.
Physio-selective smoothing is the recommended default because it matches
signal timescales: EDA/PPG are slow enough to benefit from 45s context, while
gaze, pupil, and IMU should preserve faster dynamics.
This is a useful lesson beyond GroupAffect-4.
In multimodal physiology, the temptation is to choose one temporal window for
all signals so the pipeline is simple.
Our results suggest that this simplicity can be costly: the window that makes
EDA more reliable may erase movement and oculomotor evidence.
Modality-aware preprocessing is therefore a low-complexity alternative to
larger models.

\paragraph{Temporal models need label-free sequence learning.}
Raw-sequence encoders are underpowered when trained only on labelled windows.
The unlabelled pool is therefore most useful for representation learning, not
for overwhelming the supervised objective with pseudo-labels.
This finding also reframes the role of unlabelled data.
Rather than forcing all pool windows into a supervised objective, temporal
pre-training lets the pool teach invariances while the labelled windows retain
control over class boundaries.
That separation is especially important for affective self-report, where
pseudo-labels are uncertain and person-specific.

\paragraph{Dimension-specific tuning is promising but risky.}
The best temporal configurations beat SVM+AP1 on Arousal and Dominance, but
the configurations are selected from a grid.
For a submission-ready temporal paper, all configuration choices should be
made on validation folds before final T3 testing.
The result is still scientifically useful because it identifies where temporal
models are most likely to pay off.
Arousal benefits from phasic physiological context; Dominance benefits from
slow recurrent structure and contrastive pre-training.
Valence remains hardest for temporal encoders, perhaps because the summary
features already capture the relevant gaze and pupil statistics at this scale.

\paragraph{Known-team scope.}
The canonical split is a temporal split within known teams, not a new-group
deployment test.
Strict group-isolated evaluation (LOGO-CV) should be treated as a separate
stress test before making deployment claims.

\section{Limitations and Next Steps}
\label{sec:limitations}

The current temporal results mix mature preprocessing findings with
exploratory sequence-modelling results.
The SimCLR/MSM factorial should be rerun with strict split isolation,
validation-only model selection, and clustered confidence intervals.
The preprocessing claims are more mature and can stand as a focused workshop
contribution.
Several additional limitations remain.
First, the dataset has only 10 sessions, so confidence intervals should be
clustered by group rather than by window in future work.
Second, the current temporal encoders use relatively lightweight architectures;
larger transformers may require more unlabelled data or stronger regularisation.
Third, centred smoothing is offline and should be converted to a causal filter
for real-time group support systems.
Finally, modality ablations are drop-one tests, so they measure marginal
contribution under the current feature set rather than standalone modality
performance.

\section{Conclusion}
\label{sec:conclusion}

For GroupAffect-4, the strongest practical lesson is that physiological
preprocessing must respect modality timescales.
Smoothing EDA and PPG while preserving faster gaze, pupil, and IMU features
improves all VAD dimensions and provides a strong default benchmark.
Including speech in the current small-sample setting does not improve the
headline benchmark and is therefore best treated as an optional analysis track.
Temporal self-supervision is a promising next track, especially for Arousal and
Dominance, but it should be treated as a separate sequence-modelling problem
with stricter validation.

\bibliographystyle{ACM-Reference-Format}
\bibliography{icmi2026_refs}

\end{document}
